\documentclass[11pt]{article}
\usepackage{amsmath, amssymb, amsfonts}
\usepackage{bbm}
\usepackage{physics}
\usepackage{geometry}
\usepackage{hyperref}
\geometry{margin=1in}

\title{Exercise Set: Matrices and Applications in Finance\\[0.3em]
\large Solutions}
\author{}
\date{}

\begin{document}
\maketitle

\section{Solutions for Matrix}

\subsection*{Fundamentals (1h30 in class)}

\paragraph{Exercise 1 — Basic Matrix Operations.}

We have
\[
A=\begin{pmatrix}2 & -1 \\ 3 & 4\end{pmatrix},\qquad 
B=\begin{pmatrix}1 & 0 \\ 2 & 1\end{pmatrix},\qquad 
v = \begin{pmatrix}3 \\ -1\end{pmatrix}.
\]

Basic operations:
\begin{align*}
A+B &= 
\begin{pmatrix}2+1 & -1+0 \\ 3+2 & 4+1\end{pmatrix}
=
\begin{pmatrix}3 & -1 \\ 5 & 5\end{pmatrix},\\[1ex]
A-B &= 
\begin{pmatrix}2-1 & -1-0 \\ 3-2 & 4-1\end{pmatrix}
=
\begin{pmatrix}1 & -1 \\ 1 & 3\end{pmatrix}.
\end{align*}

Matrix multiplication (note that in general $AB\neq BA$):
\begin{align*}
AB &= 
\begin{pmatrix}2 & -1 \\ 3 & 4\end{pmatrix}
\begin{pmatrix}1 & 0 \\ 2 & 1\end{pmatrix}
=
\begin{pmatrix}
2\cdot1 + (-1)\cdot2 & 2\cdot0 + (-1)\cdot1\\
3\cdot1 + 4\cdot2 & 3\cdot0 + 4\cdot1
\end{pmatrix}
=
\begin{pmatrix}0 & -1 \\ 11 & 4\end{pmatrix},\\[1ex]
BA &= 
\begin{pmatrix}1 & 0 \\ 2 & 1\end{pmatrix}
\begin{pmatrix}2 & -1 \\ 3 & 4\end{pmatrix}
=
\begin{pmatrix}
1\cdot2 + 0\cdot3 & 1\cdot(-1) + 0\cdot4\\
2\cdot2 + 1\cdot3 & 2\cdot(-1) + 1\cdot4
\end{pmatrix}
=
\begin{pmatrix}2 & -1 \\ 7 & 2\end{pmatrix}.
\end{align*}

Transposes:
\[
A^T = \begin{pmatrix}2 & 3 \\ -1 & 4\end{pmatrix},\qquad
B^T = \begin{pmatrix}1 & 2 \\ 0 & 1\end{pmatrix}.
\]

Determinant and inverse of $A$:
\[
\det(A) = 2\cdot 4 - (-1)\cdot 3 = 8+3 = 11,
\]
\[
A^{-1} = \frac1{\det(A)}
\begin{pmatrix}
4 & 1 \\ -3 & 2
\end{pmatrix}
=
\begin{pmatrix}
\frac{4}{11} & \frac{1}{11} \\
-\frac{3}{11} & \frac{2}{11}
\end{pmatrix}.
\]
A non-zero determinant means $A$ is invertible.

Action on the vector $v$:
\[
Av =
\begin{pmatrix}2 & -1 \\ 3 & 4\end{pmatrix}
\begin{pmatrix}3 \\ -1\end{pmatrix}
=
\begin{pmatrix}7 \\ 5\end{pmatrix},\qquad
Bv =
\begin{pmatrix}1 & 0 \\ 2 & 1\end{pmatrix}
\begin{pmatrix}3 \\ -1\end{pmatrix}
=
\begin{pmatrix}3 \\ 5\end{pmatrix}.
\]
So $A$ and $B$ transform the vector $v$ differently (rotation/shear/scaling).

\bigskip


\paragraph{Exercise 2 — Image, Kernel and Orthogonality.}

Given
\[
A=\begin{pmatrix}
1 & -1 & 0 \\
-1 & 2 & -1 \\
0 & -1 & 1
\end{pmatrix},
\]
we solve step by step.

\begin{enumerate}
\item \textbf{Kernel and image.}

We look for $x=(x_1,x_2,x_3)^\top$ such that $Ax=0$:
\[
\begin{pmatrix}
1 & -1 & 0 \\
-1 & 2 & -1 \\
0 & -1 & 1
\end{pmatrix}
\begin{pmatrix}
x_1\\x_2\\x_3
\end{pmatrix}
=
\begin{pmatrix}
0\\0\\0
\end{pmatrix}.
\]
This gives the system
\[
\begin{cases}
x_1 - x_2 = 0,\\
- x_1 + 2x_2 - x_3 = 0,\\
- x_2 + x_3 = 0.
\end{cases}
\]

From the first and third equations:
\[
x_1 = x_2,\qquad -x_2 + x_3 = 0 \Rightarrow x_3 = x_2.
\]

Set $x_2 = t$; then $x_1 = t$, $x_3 = t$, so
\[
x = t\begin{pmatrix}1\\1\\1\end{pmatrix},\qquad t\in\mathbb{R}.
\]

Thus
\[
\ker(A) = \operatorname{span}\left\{\begin{pmatrix}1\\1\\1\end{pmatrix}\right\}.
\]

By rank–nullity, $\dim(\ker A)=1$ implies $\operatorname{rank}(A)=3-1=2$.  
Therefore $\mathrm{Im}(A)$ is a 2-dimensional subspace of $\mathbb{R}^3$.

We will identify it in the next items.

\item \textbf{Verification of} $\mathrm{Im}(A) = (\ker A)^\perp$.

We know
\[
\ker(A) = \operatorname{span}\{(1,1,1)^\top\}
\quad\Longrightarrow\quad
(\ker A)^\perp = \{y\in\mathbb{R}^3 : y_1 + y_2 + y_3 = 0\}.
\]

We check that any column of $A$ is orthogonal to $(1,1,1)^\top$:
\[
(1,1,1)^\top A
=
\begin{pmatrix}1 & 1 & 1\end{pmatrix}
\begin{pmatrix}
1 & -1 & 0 \\
-1 & 2 & -1 \\
0 & -1 & 1
\end{pmatrix}
=
\begin{pmatrix}
0 & 0 & 0
\end{pmatrix}.
\]
So each column of $A$ has sum of components equal to $0$, hence lies in
\(\{y : y_1+y_2+y_3=0\}\).  

Therefore \(\mathrm{Im}(A) \subset (\ker A)^\perp\).  
Since both spaces have dimension 2, we actually have equality:
\[
\mathrm{Im}(A) = (\ker A)^\perp = \{y\in\mathbb{R}^3 : y_1 + y_2 + y_3 = 0\}.
\]

\item \textbf{Symmetry of $A$.}

We check directly
\[
A^\top
=
\begin{pmatrix}
1 & -1 & 0 \\
-1 & 2 & -1 \\
0 & -1 & 1
\end{pmatrix}
= A.
\]
So $A$ is symmetric: $A^\top=A$.

\item \textbf{Using the fundamental identity.}

In general, for any real matrix $M$,
\[
\mathrm{Im}(M^\top) = (\ker M)^\perp.
\]

For our matrix $A$, we have $A^\top = A$, and we already know $\ker(A)=\operatorname{span}\{(1,1,1)\}$. Therefore
\[
\mathrm{Im}(A^\top) = (\ker A)^\perp
= \{y: y_1+y_2+y_3=0\}.
\]
But $A^\top = A$, so
\[
\mathrm{Im}(A) = \mathrm{Im}(A^\top) = (\ker A)^\perp.
\]

This is consistent with what we found explicitly in (1)–(2).

\item \textbf{Geometric interpretation.}

\begin{itemize}
  \item The \emph{kernel} is the line
  \[
  \ker(A) = \{ t(1,1,1)^\top : t\in\mathbb{R}\},
  \]
  i.e.\ all vectors proportional to $(1,1,1)$. These are exactly the vectors that are sent to $0$ by $A$.

  \item The \emph{image} is the plane
  \[
  \mathrm{Im}(A) = \{y\in\mathbb{R}^3 : y_1+y_2+y_3=0\},
  \]
  which is the orthogonal complement of $(1,1,1)$. On this 2D plane, $A$ acts non-trivially (and in fact invertibly: restricted to this plane, $A$ has rank 2 and behaves like a positive definite operator).

  \item Geometrically: $A$ is like a “discrete Laplacian” that kills constant vectors (all components equal) but deforms (stretches/compresses) fluctuations around the mean on the plane of zero-sum vectors.
\end{itemize}

So $A$ annihilates the constant direction $(1,1,1)$ and acts non-trivially on its orthogonal plane of vectors with zero sum.
\end{enumerate}


\paragraph{Exercise 3 — Bootstrapping Zero-Coupon Prices via a Linear System.}

We have three maturities $T_1<T_2<T_3$ and
\[
P_k = D(T_k),\qquad
D(T_1)=y_1,\quad D(T_2)=y_1+y_2,\quad D(T_3)=y_1+y_2+y_3.
\]

\begin{enumerate}
\item[(1)] In vector form:
\[
D=
\begin{pmatrix}
D(T_1)\\D(T_2)\\D(T_3)
\end{pmatrix},
\qquad
y=
\begin{pmatrix}
y_1\\y_2\\y_3
\end{pmatrix},
\]
so
\[
D = Ly,\qquad
L =
\begin{pmatrix}
1 & 0 & 0\\
1 & 1 & 0\\
1 & 1 & 1
\end{pmatrix}.
\]
Since $P_k=D(T_k)$, we also have $P=Ly$ with $P=(P_1,P_2,P_3)^T$.

\item[(2)] The linear system is therefore
\[
\begin{pmatrix}
1 & 0 & 0\\
1 & 1 & 0\\
1 & 1 & 1
\end{pmatrix}
\begin{pmatrix}
y_1\\y_2\\y_3
\end{pmatrix}
=
\begin{pmatrix}
P_1\\P_2\\P_3
\end{pmatrix}
=
\begin{pmatrix}
0.97\\0.92\\0.855
\end{pmatrix}.
\]
$L$ is lower triangular with diagonal entries equal to $1$, so it is invertible and the solution is unique.

\item[(3)] Solving (forward substitution):
\[
y_1 = P_1 = 0.97,
\]
\[
y_1 + y_2 = P_2 \;\Rightarrow\; y_2 = P_2 - P_1 = 0.92 - 0.97 = -0.05,
\]
\[
y_1 + y_2 + y_3 = P_3 \;\Rightarrow\; y_3 = P_3 - P_2 = 0.855 - 0.92 = -0.065.
\]

\item[(4)] The discount factors are
\[
D(T_1)=y_1=0.97,\qquad
D(T_2)=y_1+y_2=0.92,\qquad
D(T_3)=y_1+y_2+y_3=0.855.
\]
Note: the negative $y_2,y_3$ reflect the way we parameterized $D(T)$ as cumulative sums; the actual discount curve $D(T_k)$ remains positive and decreasing.
\end{enumerate}

\bigskip

\paragraph{Exercise 4 — Finance-Motivated Linear Algebra.}

\begin{enumerate}
\item Two-asset portfolio:
\[
R_p = w^\top r = 0.4\,r_1 + 0.6\,r_2.
\]
This is the usual weighted average return.

\item Currency conversion:

\[
C =
\begin{pmatrix}
0.79 & 0 \\
0 & 0.86
\end{pmatrix},
\qquad
h =
\begin{pmatrix}
10\,000\\
8\,000
\end{pmatrix}.
\]
Value in GBP:
\[
Ch =
\begin{pmatrix}
0.79\cdot 10\,000\\
0.86\cdot 8\,000
\end{pmatrix}
=
\begin{pmatrix}
7\,900\\
6\,880
\end{pmatrix}.
\]

\item Covariance (unnormalized) of returns matrix $X$:
\[
X =
\begin{pmatrix}
0.01 & -0.03 \\
0.02 & 0.01 \\
-0.01 & 0
\end{pmatrix},
\qquad
X^T X =
\begin{pmatrix}
0.0006 & -0.0001\\
-0.0001 & 0.001
\end{pmatrix}.
\]
Up to a scalar factor, $X^TX$ is the empirical covariance matrix of the two return series.
\end{enumerate}

\bigskip

\paragraph{Exercise 5 — Portfolio Variance.}

\[
\Sigma =
\begin{pmatrix}
0.04 & 0.018\\
0.018 & 0.09
\end{pmatrix},
\qquad
w =
\begin{pmatrix}
0.6\\0.4
\end{pmatrix}.
\]

Variance:
\begin{align*}
w^\top \Sigma w
&= 
\begin{pmatrix}0.6 & 0.4\end{pmatrix}
\begin{pmatrix}
0.04 & 0.018\\
0.018 & 0.09
\end{pmatrix}
\begin{pmatrix}
0.6\\0.4
\end{pmatrix}\\
&= 0.6^2\cdot 0.04 + 0.4^2\cdot 0.09 + 2\cdot 0.6\cdot 0.4\cdot 0.018\\
&= 0.0144 + 0.0144 + 0.00864\\
&= 0.03744.
\end{align*}
Interpretation: the two individual variances each contribute $0.0144$, and the positive covariance adds $0.00864$ of additional risk due to imperfect diversification.

\bigskip

\paragraph{Exercise 6 — Diagonalization.}

\[
C = \begin{pmatrix}1 & 0.8\\ 0.8 & 1\end{pmatrix}.
\]

Characteristic polynomial:
\[
\det(C-\lambda I)
=
(1-\lambda)^2 - 0.64
=0.
\]
So
\[
\lambda_1 = 1.8,\qquad
\lambda_2 = 0.2.
\]

Corresponding eigenvectors (unnormalized):
\begin{itemize}
\item For $\lambda_1=1.8$, solve $(C-1.8I)v=0$:
\[
\begin{pmatrix}
-0.8 & 0.8\\
0.8 & -0.8
\end{pmatrix}
\begin{pmatrix}v_1\\v_2\end{pmatrix}
=0
\Rightarrow v_1=v_2,
\]
so e.g. $v_1=(1,1)^T$.

\item For $\lambda_2=0.2$, solve $(C-0.2I)v=0$:
\[
\begin{pmatrix}
0.8 & 0.8\\
0.8 & 0.8
\end{pmatrix}
\begin{pmatrix}v_1\\v_2\end{pmatrix}
=0
\Rightarrow v_1=-v_2,
\]
so e.g. $v_2=(1,-1)^T$.
\end{itemize}

Thus
\[
C = PDP^{-1},\qquad
D=\begin{pmatrix}1.8 & 0\\ 0 & 0.2\end{pmatrix},\quad
P=\begin{pmatrix}1 & 1\\ 1 & -1\end{pmatrix}.
\]


\bigskip

\paragraph{Exercise 7 --- PCA and Dimensionality Reduction (Simplified).}

We are given the \emph{centered} data matrix
\[
X =
\begin{pmatrix}
-1 & -1 \\
0 & 0 \\
1 & 1
\end{pmatrix},
\]
with $n=3$ observations and $d=2$ features.

\begin{enumerate}
\item Empirical covariance:
\[
\widehat{\Sigma} = \frac{1}{n-1} X^\top X = \frac12 X^\top X.
\]
Compute $X^\top X$:
\[
X^\top X =
\begin{pmatrix}
(-1)^2+0^2+1^2 & (-1)(-1)+0\cdot0+1\cdot1\\
(-1)(-1)+0\cdot0+1\cdot1 & (-1)^2+0^2+1^2
\end{pmatrix}
=
\begin{pmatrix}
2 & 2\\
2 & 2
\end{pmatrix},
\]
so
\[
\widehat{\Sigma} = \frac12
\begin{pmatrix}
2 & 2\\
2 & 2
\end{pmatrix}
=
\begin{pmatrix}
1 & 1\\
1 & 1
\end{pmatrix}.
\]

\item Eigenvalues/eigenvectors of $\widehat{\Sigma}$.

Characteristic polynomial:
\[
\det(\widehat{\Sigma}-\lambda I)
=
\det\begin{pmatrix}
1-\lambda & 1\\
1 & 1-\lambda
\end{pmatrix}
=
(1-\lambda)^2 - 1.
\]
So
\[
(1-\lambda)^2 - 1 = 0
\Rightarrow
(1-\lambda-1)(1-\lambda+1)=0
\Rightarrow
\lambda_1=2,\quad \lambda_2=0.
\]

For $\lambda_1=2$:
\[
\begin{pmatrix}
-1 & 1\\
1 & -1
\end{pmatrix}
\begin{pmatrix}v_1\\v_2\end{pmatrix}=0
\Rightarrow v_1=v_2,
\]
so e.g. $v^{(1)}=(1,1)^T$.

For $\lambda_2=0$:
\[
\begin{pmatrix}
1 & 1\\
1 & 1
\end{pmatrix}
\begin{pmatrix}v_1\\v_2\end{pmatrix}=0
\Rightarrow v_1=-v_2,
\]
so e.g. $v^{(2)}=(1,-1)^T$.

Unit principal directions:
\[
u_1 = \frac{1}{\sqrt{2}}\begin{pmatrix}1\\1\end{pmatrix},
\qquad
u_2 = \frac{1}{\sqrt{2}}\begin{pmatrix}1\\-1\end{pmatrix}.
\]

\item The first principal component is the eigenvector associated with the largest eigenvalue $\lambda_1=2$, i.e.\ $u_1$.

\item Projection onto the first principal component:

The 1D score of an observation $x_i$ is
\[
z_i = u_1^\top x_i.
\]
Compute:
\[
z_1 = \frac{1}{\sqrt{2}}(-1-1) = -\sqrt{2},\qquad
z_2 = \frac{1}{\sqrt{2}}(0+0) = 0,\qquad
z_3 = \frac{1}{\sqrt{2}}(1+1) = \sqrt{2}.
\]

So in 1D we have the points $-\sqrt{2}, 0, \sqrt{2}$.

\item Interpretation:
\begin{itemize}
\item \emph{Preserved:} almost all variance (here, all non-zero variance) is along the $u_1$ direction (both coordinates move together). Projecting onto $u_1$ keeps this main direction.
\item \emph{Lost:} any variation orthogonal to $u_1$ (here none, because data lie exactly on the line $X_1=X_2$).
\end{itemize}

\item PCA is useful here because:
\begin{itemize}
\item we can store each point with 1 coordinate instead of 2;
\item we can plot the data on a line and still see the main structure;
\item in noisier examples, directions with small variance (like $u_2$ here) often correspond to noise, which PCA discards.
\end{itemize}
\end{enumerate}

\bigskip
\paragraph{Exercise 8 — Empirical PCA with Python (with \texttt{pandas}).}

We observe $n$ realizations $x_i \in \mathbb{R}^d$ of a $d$-dimensional random vector $X$ and store them row-wise in a data matrix $X \in \mathbb{R}^{n\times d}$.

\begin{enumerate}

\item \textbf{Empirical covariance for $d=2$.}

Write $x_i=(x_{i1},x_{i2})^\top$ and
\[
\hat\mu=\frac1n\sum_{i=1}^n x_i
=
\begin{pmatrix}
\hat\mu_1\\
\hat\mu_2
\end{pmatrix},
\qquad
\hat\mu_j=\frac1n\sum_{i=1}^n x_{ij},\ j=1,2.
\]

The empirical covariance matrix is
\[
\widehat{\Sigma}
=\frac{1}{n-1}\sum_{i=1}^n
(x_i-\hat\mu)(x_i-\hat\mu)^\top
=
\begin{pmatrix}
\widehat{\mathrm{Var}}(X_1) & \widehat{\mathrm{Cov}}(X_1,X_2) \\
\widehat{\mathrm{Cov}}(X_2,X_1) & \widehat{\mathrm{Var}}(X_2)
\end{pmatrix}.
\]

For $d=2$ we have explicitly
\[
\widehat{\mathrm{Var}}(X_1)=\frac{1}{n-1}\sum_{i=1}^n (x_{i1}-\hat\mu_1)^2,
\quad
\widehat{\mathrm{Var}}(X_2)=\frac{1}{n-1}\sum_{i=1}^n (x_{i2}-\hat\mu_2)^2,
\]
\[
\widehat{\mathrm{Cov}}(X_1,X_2)
=\frac{1}{n-1}\sum_{i=1}^n (x_{i1}-\hat\mu_1)(x_{i2}-\hat\mu_2),
\]
and symmetry gives
\(\widehat{\mathrm{Cov}}(X_2,X_1)=\widehat{\mathrm{Cov}}(X_1,X_2)\).

\bigskip

\item \textbf{Python/\texttt{pandas} computation of mean and covariance.}

We now use \texttt{pandas} DataFrames instead of bare NumPy arrays.
Assume the CSV file \texttt{data.csv} contains $n$ rows, $d$ columns.

\begin{verbatim}
import pandas as pd
import numpy as np

# (a) load data into a pandas DataFrame
df = pd.read_csv("data.csv")      # each column = one variable X_j
# optional: df = pd.read_csv("data.csv", header=None)

n, d = df.shape

# (b) empirical mean vector (as pandas Series)
mu_hat = df.mean(axis=0)

# (c) center the data
X_centered = df - mu_hat

# (d) empirical covariance via matrix formula in pandas/Numpy
Sigma_hat = (X_centered.T @ X_centered) / (n - 1)

print("Empirical covariance matrix:")
print(Sigma_hat)

# (e) check using pandas built-in covariance function
Sigma_pandas = df.cov()

print("pandas DataFrame .cov():")
print(Sigma_pandas)
\end{verbatim}

Comments:
\begin{itemize}
  \item \verb|df.mean(axis=0)| computes column-wise means (one per variable $X_j$).
  \item Centering \verb|df - mu_hat| is done column by column automatically.
  \item \verb|df.cov()| directly implements the empirical covariance with the divisor $n-1$.
\end{itemize}

\bigskip

\item \textbf{Eigenvalues, eigenvectors and PCA projection (without black-box PCA).}

We now only use NumPy for the linear algebra stage.

\begin{verbatim}
# convert covariance matrix to NumPy array
Sigma = Sigma_hat.to_numpy()

# (a) eigen-decomposition of symmetric covariance matrix
eigvals, eigvecs = np.linalg.eigh(Sigma)

# (b) sort in decreasing order
idx = np.argsort(eigvals)[::-1]
eigvals_sorted = eigvals[idx]
eigvecs_sorted = eigvecs[:, idx]

# (c) first principal component
first_pc = eigvecs_sorted[:, 0]

# (d) projection onto first k principal components
k = 2
W = eigvecs_sorted[:, :k]       # d x k matrix of eigenvectors

# center again as NumPy matrix
Xc = X_centered.to_numpy()

# projected data in R^k
Z = Xc @ W
\end{verbatim}

Here
\begin{itemize}
  \item the columns of \(W\) are the first principal directions,
  \item the rows of \(Z\) are the coordinates of each observation in reduced dimension.
\end{itemize}

\bigskip

\item \textbf{Interpretation of the PCA results.}

Let the sorted eigenvalues of $\widehat{\Sigma}$ be
\[
\lambda_1 \ge \lambda_2 \ge \dots \ge \lambda_d.
\]

\begin{itemize}
  \item The proportion of variance explained by the first principal component is
  \[
  \text{Explained variance (PC1)}
  = \frac{\lambda_1}{\lambda_1+\cdots+\lambda_d}.
  \]
  In Python:
\begin{verbatim}
total_variance = eigvals_sorted.sum()
ratio_pc1 = eigvals_sorted[0] / total_variance
\end{verbatim}

  \item Reducing from dimension $d$ to $1$ or $2$ is reasonable when
  \[
  \frac{\lambda_1}{\sum_j \lambda_j}
  \quad\text{or}\quad
  \frac{\lambda_1+\lambda_2}{\sum_j \lambda_j}
  \]
  is large (e.g. above \(80\%-90\%\)).

  \item \textbf{In finance (asset returns):}
  \begin{itemize}
    \item PCA reveals dominant risk factors (market mode, sector factors),
    \item low-variance directions often correspond to noise — removing them denoises the data,
    \item plotting the first two PCs gives a 2D visualization of high-dimensional portfolios.
  \end{itemize}
\end{itemize}

Thus, using \texttt{pandas} simplifies data handling, while PCA itself relies on standard linear algebra (eigenvalues/eigenvectors of the empirical covariance matrix).
\end{enumerate}
\bigskip

\paragraph{Exercise 9 — Matrix Exponential in a Coupled Linear ODE.}

\[
\dot{x}(t) = A x(t),\quad
A = \begin{pmatrix}-0.2 & 0.3\\ 0 & -0.5\end{pmatrix},\quad
x(0) = \begin{pmatrix}0.1\\0.4\end{pmatrix}.
\]

\begin{enumerate}
\item Eigenvalues: since $A$ is upper triangular, its eigenvalues are
\[
\lambda_1=-0.2,\qquad \lambda_2=-0.5.
\]
They are distinct, so $A$ is diagonalizable.

\item Because $A$ is upper triangular with a single nonzero off-diagonal entry, we can compute $e^{At}$ explicitly (or via Putzer’s method). One finds
\[
e^{At} =
\begin{pmatrix}
e^{-0.2 t} & e^{-0.2 t} - e^{-0.5 t}\\
0 & e^{-0.5 t}
\end{pmatrix}.
\]

Thus
\[
x(t) = e^{At}x(0)
=
\begin{pmatrix}
e^{-0.2 t} & e^{-0.2 t} - e^{-0.5 t}\\
0 & e^{-0.5 t}
\end{pmatrix}
\begin{pmatrix}0.1\\0.4\end{pmatrix}
=
\begin{pmatrix}
0.5 e^{-0.2 t} - 0.4 e^{-0.5 t}\\
0.4 e^{-0.5 t}
\end{pmatrix}.
\]

\item Asymptotic behaviour: since both eigenvalues are negative, both $e^{-0.2t}$ and $e^{-0.5t}$ tend to $0$ as $t\to\infty$. Hence
\[
x(t)\to 0\quad \text{as }t\to\infty.
\]
The system is stable and trajectories decay to the origin.
\end{enumerate}

\bigskip


\subsection*{Advanced Applications}

\paragraph{Exercise 10 — Minimum-Variance Portfolio (numerical example).}
See exercise on Optimization.

\bigskip

\paragraph{Exercise 11 — PCA in a Three-Asset Universe.}

The covariance matrix is
\[
\Sigma=
\begin{pmatrix}
4 & 2 & 0\\
2 & 3 & 0\\
0 & 0 & 1
\end{pmatrix}.
\]

The third asset is independent of the first two (zero off-diagonal terms), so $\lambda_3=1$ is an eigenvalue with eigenvector $e_3=(0,0,1)^T$.

For the $2\times 2$ upper-left block
\[
B=
\begin{pmatrix}
4 & 2\\
2 & 3
\end{pmatrix},
\]
we solve
\[
\det(B-\lambda I)=0
\Rightarrow
\begin{vmatrix}
4-\lambda & 2\\
2 & 3-\lambda
\end{vmatrix}
= (4-\lambda)(3-\lambda)-4=0.
\]
This gives
\[
\lambda^2 -7\lambda+8=0
\Rightarrow
\lambda_{1,2} = \frac{7\pm\sqrt{17}}{2}.
\]

So the eigenvalues of $\Sigma$ are
\[
\lambda_1=\frac{7+\sqrt{17}}{2},\quad
\lambda_2=\frac{7-\sqrt{17}}{2},\quad
\lambda_3=1.
\]

The first principal component corresponds to the largest eigenvalue $\lambda_1$, with an eigenvector mainly supported on assets 1 and 2 (solving $(B-\lambda_1 I)v=0$ gives $v\propto(\lambda_1-3,2)^T$).  

\textbf{Interpretation:}
\begin{itemize}
\item Assets 1 and 2 share a strong common factor (they are correlated via the 2 in the off-diagonal).
\item Asset 3 is independent of the first two and contributes its own direction of risk ($\lambda_3=1$).
\item Most variance may be explained by the first principal component (associated with $\lambda_1$), which is a certain linear combination of assets 1 and 2.
\end{itemize}

\bigskip

\paragraph{Exercise 12 — Coupled Vasicek Model.}

We have
\[
\frac{dr(t)}{dt}=K(\theta-r(t)),
\qquad
K=
\begin{pmatrix}
2 & -1\\
0 & 3
\end{pmatrix},
\qquad
\theta=
\begin{pmatrix}
0.05\\
0.02
\end{pmatrix}.
\]

\begin{enumerate}

\item Rewrite in the form
\[
\frac{dr(t)}{dt}=Ar(t)+b,
\]
where
\[
A=-K=
\begin{pmatrix}
-2 & 1\\
0 & -3
\end{pmatrix},
\qquad
b=K\theta.
\]

\item The long-run mean satisfies
\[
0=Ar(\infty)+b
\quad\Rightarrow\quad
r(\infty)=-A^{-1}b=\theta,
\]
so the mean-reversion level is exactly $\theta$.

\item Since $A$ is upper triangular, its matrix exponential is
\[
e^{At}=
\begin{pmatrix}
e^{-2t} & e^{-2t}-e^{-3t}\\
0 & e^{-3t}
\end{pmatrix}.
\]

\item The solution of the inhomogeneous linear ODE is
\[
r(t)=\theta+e^{At}\bigl(r(0)-\theta\bigr),
\]
so explicitly
\[
\begin{aligned}
r_1(t)
&=0.05
+e^{-2t}\bigl(r_1(0)-0.05\bigr)
+\bigl(e^{-2t}-e^{-3t}\bigr)\bigl(r_2(0)-0.02\bigr),
\\[4pt]
r_2(t)
&=0.02
+e^{-3t}\bigl(r_2(0)-0.02\bigr).
\end{aligned}
\]

\item
\textbf{Financial interpretation:}
\begin{itemize}
\item $r_1$ and $r_2$ are short rates that mean-revert to $(0.05,0.02)$.
\item $r_2$ mean reverts faster (speed $3>2$).
\item As $t\to\infty$, $r(t)\to\theta$ and the effect of initial conditions disappears.
\end{itemize}

\item Bond price with deterministic rates:
\[
R(t)=r_1(t)+r_2(t),
\qquad
P(0,T)=\exp\!\left(-\int_0^T R(s)\,ds\right).
\]
Using the explicit formula for $r(t)$, $R(s)$ is a sum of exponentials, so the integral can be computed in closed form, giving $P(0,T)$ as an exponential of a linear combination of $1-e^{-2T}$ and $1-e^{-3T}$.
\end{enumerate}

\bigskip

\paragraph{Exercise 13 — Portfolio Theory and Risk.}

\begin{enumerate}
\item For a portfolio with weights $w$ and covariance matrix $\Sigma$, the variance is $w^\top\Sigma w$. In the exercise one obtains numerically
\[
w^\top\Sigma w = 0.0511,
\]
so the portfolio volatility is $\sqrt{0.0511}\approx 22.6\%$ (if returns are in yearly units).

\item To decorrelate returns, one diagonalizes $\Sigma$:
\[
\Sigma = Q D Q^\top,
\]
with $Q$ orthogonal and $D$ diagonal (eigenvalues).  
Defining new factors $z = Q^\top r$ gives
\[
\operatorname{Cov}(z)=Q^\top \Sigma Q = D,
\]
so the $z_i$ are uncorrelated principal components. These are useful for understanding independent sources of risk.
\end{enumerate}

\bigskip

\paragraph{Exercise 14 — Linear Factor Models.}

\begin{enumerate}
\item For cross-sectional regression $R=X\beta+\varepsilon$ where $X=[\mathbf 1\ \ F]$ contains a column of ones and factor $F$, the OLS estimator is
\[
\begin{pmatrix}
\hat\alpha\\
\hat\beta
\end{pmatrix}
=(X^\top X)^{-1}X^\top R.
\]
$\hat\alpha$ is the estimated intercept, $\hat\beta$ the factor loading.

\item Given a loading matrix $B$ and asset returns $r_t$, the factor estimate (in the least squares sense) solves
\[
\hat F_t=\arg\min_{f}\|r_t-Bf\|^2.
\]
Setting the gradient to zero:
\[
B^\top (r_t - B f)=0
\Rightarrow
(B^\top B) f = B^\top r_t
\Rightarrow
\hat F_t=(B^\top B)^{-1}B^\top r_t.
\]
\end{enumerate}

\bigskip

\paragraph{Exercise 15 — Markov Chains and Probability Viewpoint.}

Let the states be $B$ (Bull) and $R$ (Bear), and transition matrix
\[
P=
\begin{pmatrix}
p_{BB} & p_{BR}\\
p_{RB} & p_{RR}
\end{pmatrix},\qquad
p_{BB}+p_{BR}=1,\quad p_{RB}+p_{RR}=1.
\]

\begin{enumerate}
\item
\textbf{Markov property:}
\[
\mathbb P(X_{n+1}=j\mid X_n=i,X_{n-1},\dots,X_0)
=\mathbb P(X_{n+1}=j\mid X_n=i).
\]

\item
Probability of Bull tomorrow given Bull today:
\[
\mathbb P(X_1=B\mid X_0=B)=p_{BB}
\]
(e.g.\ $0.8$ in the numerical example).

\item
Two-day probability
\[
\mathbb P(X_2=B\mid X_0=B)=(P^2)_{BB},
\]
so we compute $P^2$ and read off the $(B,B)$ entry (e.g.\ $0.70$).

\item
Stationary distribution $\pi=(\pi_B,\pi_R)$ solves
\[
\pi P=\pi,\qquad \pi_B+\pi_R=1.
\]
Writing
\[
\pi_B p_{BB} + \pi_R p_{RB} = \pi_B
\]
and using $\pi_R=1-\pi_B$ gives
\[
\pi_B = \frac{p_{RB}}{p_{RB}+p_{BR}},
\qquad
\pi_R = \frac{p_{BR}}{p_{RB}+p_{BR}}.
\]
For the numeric values in the sheet one gets $\pi_B=0.6,\ \pi_R=0.4$.

\item
\textbf{Interpretation:}
$\pi$ gives the long-run fraction of time the market spends in each state, and is also the limiting distribution of $X_n$ as $n\to\infty$ (under standard irreducibility/aperiodicity assumptions).
\end{enumerate}

%%%%%%%%%%%%%%%%%%%%%%%%%%%%%%%%%%%%%%%%%%%%%%%%%%%%%%%%%%%%%%
\end{document}